\documentclass[11pt]{article}
\usepackage[margin=1in]{geometry}
\usepackage{amsmath,amssymb,amsthm,mathtools}
\usepackage{booktabs}
\usepackage{graphicx}
\usepackage{enumitem}
\usepackage[hidelinks]{hyperref}

\newtheorem{theorem}{Theorem}
\newtheorem{lemma}[theorem]{Lemma}
\newtheorem{corollary}[theorem]{Corollary}
\newtheorem{proposition}[theorem]{Proposition}
\theoremstyle{definition}
\newtheorem{definition}[theorem]{Definition}
\theoremstyle{remark}
\newtheorem{remark}[theorem]{Remark}

\newcommand{\R}{\mathbb R}
\newcommand{\E}{\mathbb E}
\newcommand{\Prob}{\mathbb P}
\newcommand{\ip}[2]{\langle #1,#2\rangle}
\newcommand{\norm}[1]{\lVert #1\rVert}
\DeclareMathOperator{\Beta}{Beta}
\DeclareMathOperator{\Law}{Law}
\newcommand{\B}{\mathrm B}
\newcommand{\Ores}{P_{\mathrm{res}}}

\title{The leading constant of the memory--sample tradeoff\\ for noiseless Gaussian regression}
\author{[AUTHOR NAME]}
\date{Draft, October 2026}

\begin{document}
\maketitle

\begin{abstract}
A streaming learner reads exact Gaussian linear equations
\(Y_t=\ip{X_t}{s}\), \(X_t\sim N(0,I_d)\), keeps at most \(M\) bits
between samples, and must output a unit vector at angular distance at
most \(\epsilon\) from \(s\). A recent preprint shows that \(M=o(d^2)\)
forces \(T\ge 2^{-16}d\log_2(1/\epsilon)\) samples. We show that the
same conclusion holds with \(2^{-16}\) replaced by any constant
\(c<2\ln 2\approx1.386\), uniformly in \(0<\epsilon\le 1/10\) (in fact with
\(\log_2(1/\epsilon)\) replaced by the slightly larger \(\log_2(1/\sin\epsilon)\)). We also
prove that randomized Kaczmarz with fixed-point registers of
\(\log_2(1/\epsilon)+\tfrac32\log_2 d+\log_2(1/\eta)+O(1)\) bits per coordinate
succeeds after \((1+\eta)\,d\ln(1/\sin^2\epsilon)\approx(1+\eta)\,2\ln2\cdot d\log_2(1/\epsilon)\) samples, for every signal.
Hence for memory in the window
\(d(\log_2(1/\epsilon)+2\log_2 d)\le M=o(d^2)\) the optimal number of
samples is \((2\ln 2+o(1))\,d\log_2(1/\epsilon)\), and randomized
Kaczmarz is first-order optimal. The improvement over the
original proof comes from three local changes: a moving split point in
its beta-mixture comparison, an exact beta-ratio estimate, and an
\(\arcsin\) form of the terminal cap bound. We also explain why the original fixed split at \(1/2\)
cannot give a constant above \(0.5176\), or above \(0.3336\) with the original
ratio estimates. We record how slowly the bound converges at finite
\(d\).
\end{abstract}

\section{Introduction}

\subsection{Model and previous bounds}
We use the model of the OpenAI preprints~\cite{OAIsub,OAImem}.

\begin{definition}[Finite-state Gaussian stream {\cite[Definition~1.1]{OAIsub}}]\label{def:model}
Fix integers \(d\ge2\), \(M\ge0\), \(T\ge0\) and \(0<\epsilon\le 1/10\).
For a signal \(s\in S^{d-1}\) the observations are \(X_t\sim N(0,I_d)\)
independent and \(Y_t=\ip{X_t}{s}\), \(1\le t\le T\). The learner reads the
pairs in order. After each sample its persistent state takes at most
\(2^M\) values. Transitions may depend on the index, on
\(d,M,T,\epsilon\), on the current state, on the \emph{entire current real pair}
\((X_t,Y_t)\), and on fresh randomness. Computation inside a
transition is unrestricted. All information retained for later
samples must be in the state. The learner stops by time \(T\) and
outputs \(\widehat s\in S^{d-1}\) as a function of the terminal state,
the stopping index and fresh randomness. Success means
\(\arccos\ip{\widehat s}{s}\le\epsilon\).
Measurability conventions are as in~\cite{OAIsub}.
\end{definition}

Throughout, \(\lambda=\log_2(1/\epsilon)\ge\log_2 10\). The model
charges bits of \emph{persistent} state. It does not charge real registers used inside one transition.
In particular a transition may read the real pair exactly, combine it with the
stored state in exact arithmetic, and store a rounded result. This is
how the upper bound in Section~\ref{sec:kaczmarz} uses the model.

The preprint~\cite{OAIsub} proves: if \(M(d)=o(d^2)\), then for all
large \(d\), every \(0<\epsilon\le1/10\) and every learner whose
success probability under a uniform signal is at least \(2/3\),
\(T\ge 2^{-16}d\log_2(1/\epsilon)\) \cite[Theorem~1.2]{OAIsub}. On
the other side, randomized Kaczmarz~\cite{Kac37,SV09} with real registers satisfies
\(\E\norm{z_t-s}^2=(1-1/d)^t\) \cite[\S1.2, eq.~(2); source \texttt{introduction.tex}:201--227]{OAImem}. This suggests
the scale \(d\ln(1/\epsilon^2)=2\ln2\cdot d\lambda\). However, the update stores
real numbers, and \cite{OAImem} states explicitly that it ``supplies no
finite-bit upper bound''. Sharan, Sidford and Valiant remark without proof that randomized Kaczmarz
``can be easily shown'' to reach error \(\epsilon\) with \(O(d\log(1/\epsilon))\) samples and \(O(d\log(1/\epsilon))\) memory
\cite[\S1.1]{SSV19}. Thus the order \(d\log(1/\epsilon)\) was known
for \(o(d^2)\) memory up to a factor of about \(2^{16}\cdot 2\ln 2\). Lemma~\ref{lem:kacz} is a written-out version
of that remark with explicit constants; we claim no novelty for it beyond the constants.

For context: with noisy labels and \(d^2/4\) bits, Sharan, Sidford and
Valiant~\cite{SSV19} proved an \(\Omega(d\log r)\) lower bound at
accuracy \(d^{-r}\), and left open whether the \(O(d\log(1/\epsilon))\) samples of randomized Kaczmarz are optimal \cite[\S1.1]{SSV19}. The order
question in their setting is addressed in~\cite{OAImem}, not here. Related
memory lower bounds for optimization are in~\cite{MSSV22,BZJ23}. The present
note does \emph{not} answer the question of~\cite{SSV19}. It sharpens
\(\Theta(d\log 1/\epsilon)\) to the leading constant \(2\ln 2\) for
\(o(d^2)\) memory. In this regime randomized Kaczmarz is first-order optimal.

\subsection{Results}

\begin{theorem}[Lower bound with the sharp constant]\label{thm:main}
Let \(M:\mathbb N\to\mathbb Z_{\ge0}\) satisfy \(M(d)=o(d^2)\), and let
\(c<2\ln2\). There is \(d_0=d_0(c,M(\cdot))\) such that for every
\(d\ge d_0\), every \(0<\epsilon\le 1/10\) and every learner of
Definition~\ref{def:model} with memory \(M(d)\), if \(S\) is uniform on
\(S^{d-1}\) independently of the rows and the learner's randomness and
\(\Prob\{\arccos\ip{\widehat s}{S}\le\epsilon\}\ge 2/3\), then
\[
 T\ \ge\ c\,d\log_2\frac1{\sin\epsilon}\ \ge\ c\,d\log_2\frac1\epsilon .
\]
\end{theorem}

\begin{lemma}[Finite-bit randomized Kaczmarz]\label{lem:kacz}
Let \(0<\eta\le1\), \(d\ge16/\eta^2\) and \(0<\epsilon\le1/10\). Put
\[
 T=\bigl\lceil(1+\eta)\,d\ln(1/\sin^2\epsilon)\bigr\rceil,\qquad
 q=\Bigl\lceil\log_2\frac1{\sin\epsilon}+\frac32\log_2 d+\log_2\frac{4\sqrt6}{\eta}\Bigr\rceil .
\]
There is a learner of Definition~\ref{def:model} with
\[
 M=d(q+3)\ \le\ d\Bigl(\lambda+\frac32\log_2 d+\log_2\frac1\eta+7.3\Bigr)
\]
bits, horizon \(T\), and a deterministic transition rule (fixed-point
randomized Kaczmarz). For \emph{every} signal \(s\in S^{d-1}\) it
satisfies \(\Prob_s\{\arccos\ip{\widehat s}{s}\le\epsilon\}\ge 2/3\).
Moreover
\[
 d\ln\frac1{\sin^2\epsilon}\ \le\ d\ln\frac1{\epsilon^2}+\frac{d\epsilon^2}{2}
 \ \le\ 1.0011\cdot 2\ln 2\cdot d\lambda .
\]
Hence \(T\le(1+\eta)\,1.0011\cdot 2\ln2\cdot d\lambda+1\).
\end{lemma}

In the form requested by the lower bound, Lemma~\ref{lem:kacz} gives
\(M\le d(\lambda+C\log_2 d)\) with
\(C=\tfrac32+(\log_2(1/\eta)+7.3)/\log_2 d\). For example, \(C=2\)
works once \(d\ge 2^{14.6}\eta^{-2}\). The range \(0<\epsilon\le 1/10\) is
uniform, and \(d_0=16/\eta^2\) does not depend on \(\epsilon\).

\begin{corollary}[Sharpness window]\label{cor:sharp}
Let \(M(d)=o(d^2)\), and let \(\epsilon=\epsilon_d\in(0,1/10]\) satisfy
\(d(\lambda_d+2\log_2 d)\le M(d)\) for all large \(d\). Let
\(T^*(d)\) be the least horizon for which some learner with memory
\(M(d)\) reaches accuracy \(\epsilon_d\) with probability \(\ge2/3\) for
every signal. (The same conclusions hold for the least horizon with success
probability \(\ge2/3\) under a uniform signal, which lies between the two bounds below.) Then
\[
 \lim_{d\to\infty}\frac{T^*(d)}{d\ln(1/\sin^2\epsilon_d)}=1,
\]
hence
\[
 2\ln2\le\liminf_{d}\frac{T^*(d)}{d\log_2(1/\epsilon_d)}\le\limsup_d\frac{T^*(d)}{d\log_2(1/\epsilon_d)}\le1.0011\cdot2\ln2,
\]
with \(\lim T^*(d)/(d\log_2(1/\epsilon_d))=2\ln2\) whenever \(\epsilon_d\to0\).
\end{corollary}
\begin{proof}
Note \(d\ln(1/\sin^2\epsilon)=2\ln2\cdot d\log_2(1/\sin\epsilon)\).
For the lower bound, an every-signal guarantee implies the uniform-prior guarantee
\cite[Corollary~1.3]{OAIsub}, so Theorem~\ref{thm:main} gives
\(\liminf T^*/(d\ln(1/\sin^2\epsilon_d))\ge c/(2\ln2)\) for every \(c<2\ln 2\).
For the upper bound, fix \(\eta\in(0,1]\). Once \(\tfrac12\log_2d\ge\log_2(1/\eta)+7.3\) and
\(d\ge16/\eta^2\), Lemma~\ref{lem:kacz} applies within memory
\(d(\lambda+2\log_2 d)\le M(d)\), giving \(T^*\le(1+\eta)d\ln(1/\sin^2\epsilon_d)+1\). Let \(\eta\downarrow0\).
The second display follows from
\(1\le\ln(1/\sin^2\epsilon)/\ln(1/\epsilon^2)\le1+\epsilon^2/(2\ln(1/\epsilon^2))\le1.0011\) (Lemma~\ref{lem:kacz}).
\end{proof}

\begin{remark}[Window and the \(\Theta(d^2)\) regime]
Corollary~\ref{cor:sharp} needs both ends of the window. If
\(M<d\lambda\), the learner cannot store \(\lambda\) bits per
coordinate, and our upper bound does not apply. If \(M=Ad^2\) with fixed
\(A>0\), Theorem~\ref{thm:main} is not available. The proof below then
yields only a smaller constant: \(c(A)\approx0.47,0.72,0.93,1.19,1.31\)
for \(A=10^{-2},10^{-3},10^{-4},10^{-6},10^{-8}\)
(script \texttt{computations/memratio.py}, asymptotic in \(d\); heuristic, since it treats the dimension loss \(2J\) as
negligible, while for \(M=\Theta(d^2)\) the short-run lemma only gives \(\lambda=O(d)\)). The sharp constant for
\(\Theta(d^2)\) memory is unknown.
\end{remark}

\subsection{What is new}
The overall architecture of the lower bound is that of~\cite{OAIsub}:
residual subspheres, an \(m\)-th moment to pay for state selection,
a two-component beta mixture for correlated replica directions,
backward induction over blocks, and a terminal cap. We change three
estimates and re-optimize the parameters.
\begin{enumerate}[label=(\roman*)]
\item \emph{Moving split} (Lemma~\ref{lem:moving}). The original
two-component domination splits the beta density at \(q_0=1/2\). This
costs a factor \(2^{(b-1)/2}\) per replica direction, which is \((\kappa/2)\ln2\)
nats per unit dimension per block, independently of all parameters. This is at least as large as the
residual cost \((\kappa/2)\ln\frac1{1-\pi}\) whenever \(\pi\le1/2\), so no choice of parameters removes it. Splitting at \(q_0=1-\delta\)
with \(\delta\to0\) reduces this cost to \((\kappa/2)\ln\frac1{1-\delta}=O(\kappa\delta)\).
\item \emph{Beta-ratio bound} (Lemma~\ref{lem:betaratio}). The price of
the moving split is a factor \(\delta^{-(b-2)/2}\) on the second
component. It is cancelled by the exact ratio
\(\B(\alpha,\beta_1)/\B(\alpha,\beta_0)\le((b+j-2)/(D-2))^{\lfloor(b-2)/2\rfloor}\).
The original proof bounded this ratio by \(1\).
\item \emph{\(\arcsin\) terminal cap} (Lemma~\ref{lem:cap}). The terminal bound
improves from \((8\epsilon/r)^{D-1}\) to \(\frac12(\sin\epsilon/r)^{D-1}\). This makes the
argument uniform down to \(\lambda=\log_2 10\) and removes the separate
small-\(\lambda\) branch of~\cite{OAIsub}.
\end{enumerate}
The original method fails at a definite step. With the split fixed at
\(1/2\), every choice of parameters gives at most
\(0.5176\), and at most \(0.3336\) with the original ratio estimates
(Section~\ref{sec:framework}). The moving split is therefore necessary,
not a re-tuning. Section~\ref{sec:ledger} records which change
contributes which factor of the total \(2^{-16}\to1.3847\)
(certified) improvement, a factor of about \(9.07\times10^4\).

\section{Three estimates}\label{sec:lemmas}

Write \(f_{A,B}\) for the \(\Beta(A,B)\) density and \(\B\) for the beta
function. Recall that \(\B(A,B)\) decreases in each argument and that
\(\B(A,B+1)=\B(A,B)\,B/(A+B)\).

\subsection{The terminal cap}

\begin{lemma}[\(\arcsin\) cap]\label{lem:cap}
Let \(H\) be a Euclidean space of dimension \(D\ge2\), let
\(w\in H\) and \(0<\rho<1\). Then
\[
 \sigma_H\{u\in S(H):\norm{u-w}\le\rho\}\ \le\ \tfrac12\,I_{\rho^2}\Bigl(\tfrac{D-1}{2},\tfrac12\Bigr)\ \le\ \tfrac12\rho^{D-1},
\]
where \(I_x(a,b)\) is the regularized incomplete beta function.
Consequently, on an affine subsphere \(z+rS(H)\) with \(z\perp H\),
\(\norm z^2+r^2=1\), the angular success probability \(\sigma_H\{u:\arccos\ip{\widehat s}{z+ru}\le\epsilon\}\)
of any fixed unit output \(\widehat s\), or of any target-independent output law, is
at most \(\min\{1,\frac12(\sin\epsilon/r)^{D-1}\}\le(\sin\epsilon/r)^p\)
for every \(0\le p\le D-1\).
\end{lemma}
\begin{proof}
If \(\norm{u-w}\le\rho<1=\norm u\), then
\(1-2\ip uw+\norm w^2\le\rho^2<1\), so \(\ip uw>\norm w^2/2\ge0\) and
\(w\ne0\). The distance from \(u\) to the line \(\R w\) is at most
\(\norm{u-w}\le\rho\). Since \(\ip uw>0\), the angle \(\phi\) between \(u\)
and \(e=w/\norm w\) is acute with \(\sin\phi\le\rho\), so
\(\phi\le\theta:=\arcsin\rho\). The cap \(\{\phi\le\theta\}\) with
\(\theta\le\pi/2\) has mass
\(\frac12\Prob(1-\ip ue^2\le\rho^2)=\frac12 I_{\rho^2}(\frac{D-1}2,\frac12)\),
because \(1-\ip ue^2\sim\Beta(\frac{D-1}2,\frac12)\). With \(a=(D-1)/2\)
and the substitution \(t=\rho^2 s\),
\[
 I_{\rho^2}(a,\tfrac12)=\frac{\rho^{2a}}{\B(a,\frac12)}\int_0^1 s^{a-1}(1-\rho^2s)^{-1/2}ds
 \le\frac{\rho^{2a}}{\B(a,\frac12)}\int_0^1 s^{a-1}(1-s)^{-1/2}ds=\rho^{D-1}.
\]
For the consequence: if \(x\in S^{d-1}\) and \(\angle(x,\widehat s)\le\epsilon\), then
\(\norm{x-\cos\epsilon\,\widehat s}^2=1-2\cos\epsilon\ip x{\widehat s}+\cos^2\epsilon\le\sin^2\epsilon\).
Apply this to the unit vector \(x=z+ru\) and project onto \(H\) (recall \(z\perp H\)):
\(\norm{u-w}\le\sin\epsilon/r\) with \(w=\cos\epsilon\,P_H\widehat s/r\). If \(\sin\epsilon/r<1\), apply the
first part with \(\rho=\sin\epsilon/r\). Otherwise use total
mass \(1\le(\sin\epsilon/r)^p\). Averaging over an output law independent of the target
preserves the bound.
\end{proof}
This replaces \cite[Lemma~3.1, eq.~(12) and its proof, \texttt{02-subspace-moments.tex}:24--56]{OAIsub}. There
the successful set was placed in a chordal ball of radius \(\epsilon/r\), recentred at a successful direction (doubling the radius), and then
bounded by the angular cap \((2\theta)^{D-1}\) with \(\theta\le(\pi/2)\cdot 2\epsilon/r\). This gives
\((2\pi\epsilon/r)^{D-1}\), rounded to \((8\epsilon/r)^{D-1}\). Using the centre \(\cos\epsilon\,\widehat s\) instead of \(\widehat s\)
also replaces \(\epsilon\) by \(\sin\epsilon\), which makes the lower bound match the Kaczmarz bound exactly at the level of
\(d\ln(1/\sin^2\epsilon)\). The worst-case centre
for a chordal ball is not on the sphere. It is at distance \(\sqrt{1-\rho^2}\) from the origin,
which is why the exact answer is the \(\arcsin\) cap. This was checked by Monte Carlo
(\texttt{computations/independent/indep\_cap\_kacz.py}, part f2).

\subsection{The moving split}
We use the setting of \cite[Lemma~5.1]{OAIsub}. Let \(E\) have dimension \(D'\), let
\(V\subset E\) have dimension \(1\le j<n<D'\), let \(N\) be uniform among
the \(n\)-dimensional subspaces of \(E\) containing \(V\), and let \(v\) be uniform
on \(S(N)\) given \(N\). Put \(b=D'+1-n\ge2\) and
\[
 \alpha=\tfrac{n-j}2,\quad \alpha_*=\tfrac{D'-j}2=\alpha+\tfrac{b-1}2,\quad
 \beta_0=\tfrac j2,\quad \beta_1=\tfrac{D'-n+j-1}2=\beta_0+\tfrac{b-2}2 .
\]
Let \(\nu_0=\sigma_E\), and let \(\nu_1\) be the law of
\(\sqrt{1-Z}\,e_0+\sqrt Z\,u_0\) with \(Z\sim\Beta(\alpha,\alpha_*-\alpha)\),
\(e_0\sim\sigma_V\), and \(u_0\sim\sigma_{E\cap e_0^\perp}\) given \(e_0\), all independent.

\begin{lemma}[Moving-split domination]\label{lem:moving}
For every \(0<\delta<1\), as measures on \(S(E)\),
\[
 \Law(v)\ \le\ A_\delta\,\nu_0+B_\delta\,\nu_1,\qquad
 A_\delta=\frac{\B(\alpha_*,\beta_0)}{\B(\alpha,\beta_0)}(1-\delta)^{-\frac{b-1}2}\le(1-\delta)^{-\frac{b-1}2},\quad
 B_\delta=\frac{\B(\alpha,\beta_1)}{\B(\alpha,\beta_0)}\,\delta^{-\frac{b-2}2}.
\]
\end{lemma}
\begin{proof}
Let \(Q(y)=\norm{P_{V^\perp}y}^2\). As in the proof of \cite[Lemma~5.1]{OAIsub}
(\texttt{07-beta-mixture.tex}:50--67), \(Q(v)\sim\Beta(\alpha,\beta_0)\),
\(Q\sim\Beta(\alpha_*,\beta_0)\) under \(\nu_0\), and \(Q\sim\Beta(\alpha,\beta_1)\)
under \(\nu_1\). The three laws are invariant under
\(O(V)\times O(V^\perp)\), and their two component lengths are a.s.\ nonzero. Hence a
pointwise inequality between the scalar densities of \(Q\) lifts to the vector
measures (\emph{ibid.}, lines 69--79). Put \(q_0=1-\delta\). For \(q\ge q_0\),
\[
 f_{\alpha,\beta_0}(q)=\frac{\B(\alpha_*,\beta_0)}{\B(\alpha,\beta_0)}\,q^{-\frac{b-1}2}f_{\alpha_*,\beta_0}(q)
 \le A_\delta\,f_{\alpha_*,\beta_0}(q),
\]
and \(\B(\alpha_*,\beta_0)\le\B(\alpha,\beta_0)\) since \(\alpha_*>\alpha\). For \(q<q_0\) we have \(1-q>\delta\), so
\[
 f_{\alpha,\beta_0}(q)=\frac{\B(\alpha,\beta_1)}{\B(\alpha,\beta_0)}(1-q)^{-\frac{b-2}2}f_{\alpha,\beta_1}(q)
 \le B_\delta\,f_{\alpha,\beta_1}(q).
\]
Hence \(f_{\alpha,\beta_0}\le A_\delta f_{\alpha_*,\beta_0}+B_\delta f_{\alpha,\beta_1}\) on \((0,1)\).
\end{proof}
At \(\delta=1/2\) and with \(\B(\alpha,\beta_1)\le\B(\alpha,\beta_0)\) this is the original
\(\Law(v)\le 2^{(b-1)/2}(\nu_0+\nu_1)\). The statement there is the weaker \(2^{D'}(\nu_0+\nu_1)\).
Monotonicity of the beta function is the only input, and it holds for every \(q_0\).

\subsection{Beta ratio and inverse moment}

\begin{lemma}[Beta-ratio bound]\label{lem:betaratio}
For \(\alpha,\beta_0>0\) and \(h\ge0\),
\(\displaystyle\frac{\B(\alpha,\beta_0+h)}{\B(\alpha,\beta_0)}\le\Bigl(\frac{\beta_0+h}{\alpha+\beta_0+h}\Bigr)^{\lfloor h\rfloor}\).
In the notation above, \(\alpha+\beta_1=(D-2)/2\) with \(D=D'+1\), so
\[
 \frac{\B(\alpha,\beta_1)}{\B(\alpha,\beta_0)}\le\Bigl(\frac{b+j-2}{D-2}\Bigr)^{\lfloor (b-2)/2\rfloor}.
\]
\end{lemma}
\begin{proof}
\(\B(\alpha,\beta_0+h)\le\B(\alpha,\beta_0+\lfloor h\rfloor)=\B(\alpha,\beta_0)\prod_{i=0}^{\lfloor h\rfloor-1}\frac{\beta_0+i}{\alpha+\beta_0+i}\).
Each factor is at most \((\beta_0+h)/(\alpha+\beta_0+h)\), since \(x\mapsto x/(\alpha+x)\) increases.
For the second claim, \(\alpha+\beta_1=\frac{n-j}2+\frac{D'-n+j-1}2=\frac{D'-1}2\) and \(\beta_1=\frac{b+j-2}2\).
\end{proof}

\begin{lemma}[Inverse moment]\label{lem:invmom}
If \(Z\sim\Beta(\alpha,h_1)\) and \(\ell<\alpha-1\) is a nonnegative integer, then
\(\E Z^{-\ell}=\prod_{t=1}^{\ell}\frac{\alpha+h_1-t}{\alpha-t}\le\bigl(\frac{\alpha-1}{\alpha-1-\ell}\bigr)^{h_1}\).
(The tempting bound \((\alpha/(\alpha-\ell))^{h_1}\) is false for \(h_1=\tfrac12\), e.g.\ \(\alpha=\tfrac52,\ell=1\): \(4/3>\sqrt{5/3}\).)
\end{lemma}
\begin{proof}
The product formula is \cite[eq.~(6)]{OAIsub}. We have
\(\ln\E Z^{-\ell}=\sum_{t=1}^\ell\ln(1+\frac{h_1}{\alpha-t})\le h_1\sum_{t=1}^\ell\frac1{\alpha-t}\)
and \(\frac1{\alpha-t}\le\int_{t}^{t+1}\frac{dx}{\alpha-x}\), since the integrand increases on \([t,t+1]\subset[0,\alpha-1)\); summing over \(1\le t\le\ell\) gives \(h_1\ln\frac{\alpha-1}{\alpha-1-\ell}\).
\end{proof}

\begin{proposition}[One replica direction]\label{prop:dir}
In the setting of Lemma~\ref{lem:moving}, let \(\ell<\alpha-1\) be an integer, and let
\(r_1\ge(b+j-2)/(D-2)\) and \(r_2\ge(\alpha-1)/(\alpha-1-\ell)\) with \(r_1<1\le r_2\).
If \(r_1r_2\le\delta<1\), then
\[
 A_\delta+B_\delta\,\E Z^{-\ell}\ \le\ (1-\delta)^{-\frac{b-1}2}+r_2\delta^{-1/2}.
\]
\end{proposition}
\begin{proof}
By Lemmas~\ref{lem:betaratio} and~\ref{lem:invmom} with \(h_1=(b-1)/2\),
\(B_\delta\E Z^{-\ell}\le r_1^{\lfloor (b-2)/2\rfloor}\delta^{-(b-2)/2}r_2^{(b-1)/2}\).
For even \(b\) this equals \((r_1r_2/\delta)^{(b-2)/2}r_2^{1/2}\le r_2^{1/2}\).
For odd \(b\ge3\) it equals \((r_1r_2/\delta)^{(b-3)/2}\delta^{-1/2}r_2\le r_2\delta^{-1/2}\).
\end{proof}

The point is the scaling. For a block of \(b\le k=\kappa d\) rows, the first term
costs \(\frac{\kappa}{2}\ln\frac1{1-\delta}\) nats per unit dimension. The second term
is \(O(1)\) per direction, hence \(O(1/d)\) per unit dimension, provided \(\delta\) exceeds
\(r_1r_2\approx(\kappa+\mu)/(1-\pi/(1-\kappa-\mu))\). Since this threshold for \(\delta\) tends to \(0\) with \(\kappa,\mu\), the first term
is \(o(\kappa)\) as \(\kappa,\mu\to0\), while the main residual cost is of order \(\kappa\pi\).

\section{Proof of Theorem~\ref{thm:main}}\label{sec:proof}

\paragraph{Rate function.} For parameters \(\kappa,\pi,\mu,\zeta>0\) with \(\kappa+\mu+\pi<1\) put
\[
 \delta=\frac{(1+\zeta)(\kappa+\mu)}{1-\pi/(1-\kappa-\mu)},\qquad
 R_1(\kappa,\pi)=\int_0^{\pi/2}\ln\frac{1/2-x}{(1-\kappa)/2-x}\,dx,
\]
\[
 g=R_1(\kappa,\pi)+\frac\kappa2\ln\frac1{1-\delta},\qquad
 \Phi(\kappa,\pi,\mu,\zeta)=\frac{\kappa\pi\ln2}{g}-\frac{\kappa}{\log_2 10},
\]
and call the parameters admissible if moreover \(\delta<1\).

\begin{lemma}\label{lem:sup}
\(\sup\{\Phi:\text{admissible}\}=2\ln 2\). Explicitly, \(\Phi(10^{-7},10^{-4},10^{-8},10^{-3})\ge1.38470\).
This was certified with 200-bit interval arithmetic (\texttt{computations/certify.py}) and recomputed independently
with \texttt{mpmath} quadrature (\texttt{computations/independent/indep\_constant.py}).
\end{lemma}
\begin{proof}
Using \(\ln(1+y)\le y\),
\(R_1\le\frac\kappa2\int_0^{\pi/2}\frac{dx}{(1-\kappa)/2-x}=\frac\kappa2\ln\frac{1-\kappa}{1-\kappa-\pi}\).
Take \(\zeta=1\), \(\mu=\kappa^2\), and fix \(\pi\in(0,\frac12)\). As \(\kappa\to0\) we have
\(\delta=O(\kappa)\), so \(g\le\frac\kappa2[\ln\frac1{1-\pi}+O(\kappa)]\) and
\(\Phi\ge\frac{2\pi\ln 2}{\ln(1/(1-\pi))}-O(\kappa)\). The right side tends to \(2\ln2\) as
\(\pi\to0\). Conversely, \(R_1\ge\frac\kappa2\ln\frac1{1-\pi}\) gives \(\Phi<2\pi\ln2/\ln\frac1{1-\pi}<2\ln 2\).
\end{proof}

\paragraph{Setup.} Write \(\lambda_s=\log_2(1/\sin\epsilon)\). Then \(\lambda\le\lambda_s\le\lambda+0.0025\) and
\(\lambda_s\ge\log_210\) for \(\epsilon\le1/10\).
Suppose the theorem fails for some \(c<2\ln2\). Then there is an unbounded
sequence of dimensions \(d\), accuracies \(\epsilon_d\le1/10\), and learners with memory \(M(d)\) and success \(\ge2/3\) such that
\(T<cd\lambda_s\le2cd\lambda\). Fix admissible \((\kappa,\pi,\mu,\zeta)\) with \(\Phi>c\) (Lemma~\ref{lem:sup}). By the
short-run lemma \cite[Lemma~4.2]{OAIsub}, applied with \(2c\) and \(\alpha=2/3\), we have
\(\lambda=o(d)\), \(T=o(d^2)\), and \(\log_2 w=o(d^2)\) for the padded width \(w=(T+2)2^M\). These bounds
are uniform in the accuracy and the rules, because \cite[eq.~(18)]{OAIsub} bounds \((d-2)\lambda\) by a
function of \(d\) and \(M(d)\). Set
\[
 k=\lfloor\kappa d\rfloor,\quad m=\lceil\mu d\rceil,\quad \ell=\lfloor\pi d/2\rfloor,\quad p=2\ell,\quad
 J=\lceil T/k\rceil\le 2c\lambda_s/\kappa+1=o(d),\quad D_{\min}=d-2J .
\]
Put \(\alpha_{\min}=(D_{\min}-k-m)/2\), \(r_1=(k+m)/(D_{\min}-2)\) and \(r_2=(\alpha_{\min}-1)/(\alpha_{\min}-1-\ell)\).
Since \(J=o(d)\), \(r_1\to\kappa+\mu\) and \(r_2\to(1-\pi/(1-\kappa-\mu))^{-1}\), so \(r_1r_2\to\delta/(1+\zeta)<\delta\).
Hence for all large \(d\) in the sequence:
\begin{equation}\label{eq:eventual}
 r_1r_2\le\delta,\qquad p\le D_{\min}-1,\qquad m<D_{\min}-k,\qquad \ell+1<\alpha_{\min}.
\end{equation}
Define
\[
 \Ores=\prod_{t=1}^{\ell}\frac{D_{\min}/2-t}{(D_{\min}-k)/2-t},\qquad
 G=w^{1/m}\bigl[(1-\delta)^{-k/2}+r_2\delta^{-1/2}\bigr]\,\Ores .
\]

\paragraph{Invariant.} Fix a seed as in \cite[Lemma~2.1]{OAIsub}, pad to horizon \(T\)
\cite[Lemma~4.1]{OAIsub}, and cut the horizon into \(J\) batches of length \(\le k\). For a state \(v\) after \(i\)
batches let \(h_{i,v}(s)\) be its suffix success. We claim that every affine subsphere \(c_0+rS(H)\) with
\(\dim H=D\ge d-2i\) satisfies
\begin{equation}\label{eq:inv}
 \int h_{i,v}(c_0+ru)\,d\sigma_H(u)\le G^{J-i}(\sin\epsilon/r)^p .
\end{equation}
At \(i=J\) this is Lemma~\ref{lem:cap}, since \(D\ge D_{\min}\ge p+1\).

\paragraph{One block.} Assume~\eqref{eq:inv} at \(i+1\). Let the batch have length \(1\le b\le k\) and let \(D\ge d-2i\).
By \cite[Lemmas~2.4 and~2.5]{OAIsub}, \(\rho^2\sim\Beta((D-b)/2,b/2)\), \(r'=r\rho\), \(E=H\cap e^\perp\), \(D'=D-1\),
\(n=D-b\). The left side of~\eqref{eq:inv} is at most \(\E_{\rho,e}\E_N\max_{v'}F_{v'}(N)\), where
\(F_{v'}(N)=\int h_{i+1,v'}(c'+r'u)\,d\sigma_N(u)\). Put \(U_0=G^{J-i-1}(\sin\epsilon/r')^p\). Expand
\(F_{v'}(N)^m\) over \(m\) conditionally independent directions. For \(j=0\) or \(b=1\) the next direction is
uniform on \(S(E)\) \cite[Lemma~3.2]{OAIsub}, and its factor is \(\le U_0\). For \(j\ge1\) and \(b\ge2\), apply
Lemma~\ref{lem:moving}. Under \(\nu_0\) the factor is \(\le U_0\), by the stage-\((i+1)\) bound in dimension
\(D-1\). Under \(\nu_1\), given \((Z,e_0)\), the target is uniform on an admissible affine subsphere of dimension
\(D-2\ge d-2(i+1)\) and radius \(r'\sqrt Z\) (\cite{OAIsub}, \texttt{07-beta-mixture.tex}:224--237), so the
factor is \(\le U_0Z^{-\ell}\). Now \(\alpha=(D-b-j)/2>\alpha_{\min}\) and \((b+j-2)/(D-2)\le r_1\), so by
Proposition~\ref{prop:dir} and~\eqref{eq:eventual} the conditional factor is at most
\([(1-\delta)^{-k/2}+r_2\delta^{-1/2}]\,U_0\). Integrating in reverse conditional order, applying Jensen's
inequality and summing over at most \(w\) states \cite[Lemma~2.5]{OAIsub} gives
\(\E_N\max_{v'}F_{v'}\le w^{1/m}[(1-\delta)^{-k/2}+r_2\delta^{-1/2}]\,U_0\).
Finally \(U_0\) contains \(\rho^{-p}\), and
\(\E\rho^{-p}=\prod_{t=1}^\ell\frac{D/2-t}{(D-b)/2-t}\le\Ores\). Each factor increases in \(b\) and decreases in
\(D\ge D_{\min}\). This proves~\eqref{eq:inv} at \(i\).

\paragraph{Conclusion.} At \(i=0\) (with \(H=\R^d\), \(r=1\)), averaging over the seed and the initial state as in \cite{OAIsub}
gives \(2/3\le G^J(\sin\epsilon)^p=G^J2^{-p\lambda_s}\). Since \(J\le T/k+1\),
\begin{equation}\label{eq:TLB}
 T\ \ge\ k\Bigl(\frac{p\lambda_s\ln2-\ln(3/2)}{\ln G}-1\Bigr).
\end{equation}
We have \(\ln G\le\frac{\ln w}{m}+\frac k2\ln\frac1{1-\delta}+\ln(1+r_2\delta^{-1/2})+\ln\Ores\). Here
\(\ln w/(md)\to0\), because \(M=o(d^2)\) and \(m\ge\mu d\); this is the only place where \(o(d^2)\) memory enters
besides the short-run lemma. The factors of \(\Ores\) increase in \(t\), so
\(\ln\Ores\le\int_0^{\ell+1}\ln\frac{D_{\min}/2-x}{(D_{\min}-k)/2-x}dx=(1+o(1))\,dR_1(\kappa,\pi)\), because \(D_{\min}/d\to1\). Thus
\(\limsup\ln G/d\le g\). Also \(\ln G\ge\ln\Ores\ge\ell\ln(1+k/d)\), which grows linearly in \(d\). Dividing~\eqref{eq:TLB}
by \(d\lambda_s\) and using \(\lambda_s\ge\log_210\) gives
\(\liminf T/(d\lambda_s)\ge\kappa\pi\ln2/g-\kappa/\log_2 10=\Phi>c\). This contradicts \(T<cd\lambda_s\).
Every eventual condition depends only on \(d\), \(M(d)\) and the fixed parameters, which gives the
threshold \(d_0(c,M)\). \qed

\begin{remark}
The small-accuracy branch of \cite{OAIsub} (\(\lambda\le1024\), handled by the full-data bound \(T>d/4\)) is no longer needed.
With terminal constant \(1\), the induction itself is effective for all \(\lambda_s\ge\log_210\). The \(-\kappa/\log_210\) term in \(\Phi\)
is the price of this uniformity, and it vanishes as \(\kappa\to0\).
\end{remark}

\section{Finite-bit Kaczmarz: proof of Lemma~\ref{lem:kacz}}\label{sec:kaczmarz}

\paragraph{The learner.} Let \(\mathcal G_q=2^{-q}\mathbb Z\cap[-3,3]\). Then
\(|\mathcal G_q|=6\cdot2^q+1\le2^{q+3}\), and the state space is \(\mathcal G_q^d\), so \(M=d(q+3)\) bits suffice.
Start with \(z_0=0\). On reading the real pair \((x_t,y_t)\) with \(x_t\ne0\), compute in exact arithmetic
\[
 \tilde z_t=z_{t-1}+\frac{y_t-\ip{x_t}{z_{t-1}}}{\norm{x_t}^2}x_t ,
\]
and store \(z_t=\mathrm{rd}_q(\tilde z_t)\), where \(\mathrm{rd}_q\) rounds each coordinate to the nearest point of \(2^{-q}\mathbb Z\)
(ties toward \(0\); a Borel rule). On the null event \(x_t=0\) keep \(z_t=z_{t-1}\). After \(T\) samples output
\(z_T/\norm{z_T}\), or a fixed unit vector if \(z_T=0\). This is a learner of Definition~\ref{def:model}: the real sample is
used only inside the transition, and only \(z_t\) is retained. We show below that \(\tilde z_t\in[-3,3]^d\), so \(z_t\in\mathcal G_q^d\).

\paragraph{Error recursion.} Fix \(s\in S^{d-1}\) and put \(u_t=x_t/\norm{x_t}\), \(P_t=I-u_tu_t^{\mathsf T}\),
\(e_t=z_t-s\). Since \(y_t=\ip{x_t}s\) exactly, \(\tilde z_t-s=P_te_{t-1}\). Hence
\[
 e_t=P_te_{t-1}+r_t,\qquad \norm{r_t}_\infty\le2^{-q-1},\qquad \norm{r_t}_2\le\delta_q:=\sqrt d\,2^{-q-1}.
\]
Let \(\hat e_t=P_t\cdots P_1(-s)\) be the error of exact Kaczmarz on the same rows, and let \(a_t=e_t-\hat e_t\). Then
\(a_t=P_ta_{t-1}+r_t\), \(a_0=0\), so
\(a_T=\sum_{t=1}^TP_T\cdots P_{t+1}r_t\). As \(\norm{P_t}\le1\), we get \(\norm{a_t}\le t\delta_q\) deterministically.

\paragraph{Range.} We have \(\norm{e_{t-1}}\le\norm{\hat e_{t-1}}+\norm{a_{t-1}}\le1+T\delta_q\), and so
\(\norm{\tilde z_t}_\infty\le\norm s+\norm{P_te_{t-1}}\le2+T\delta_q\). We claim \(T\delta_q\le1\). Indeed
\(T\le(1+\eta)d\Lambda+1\) with \(\Lambda=\ln(1/\sin^2\epsilon)=2\ln2\cdot\lambda'\) and \(\lambda'=\log_2(1/\sin\epsilon)\). Hence
\(\log_2(T\sqrt d/2)\le\frac32\log_2d+\log_2(2.78\lambda'+1)-1\). Since \(2^{\lambda'+4.29}\ge19.5(1+0.69\lambda')\ge2.78\lambda'+1\), this is
\(\le\frac32\log_2d+\lambda'+3.29\le q\). So \(\tilde z_t\in[-3,3]^d\), and rounding to the nearest grid point stays in
\([-3,3]\) because \(\pm3\in\mathcal G_q\).

\paragraph{Exact part: a Beta product.} Let \(\mathcal F_t=\sigma(x_1,\dots,x_t)\). If \(\hat e_{t-1}\ne0\), put
\(B_t=\ip{u_t}{\hat e_{t-1}/\norm{\hat e_{t-1}}}^2\). Otherwise put \(B_t=\ip{u_t}{e_1}^2\). Since \(u_t\) is uniform and
independent of \(\mathcal F_{t-1}\), \(B_t\sim\Beta(\frac12,\frac{d-1}2)\) independently of \(\mathcal F_{t-1}\). Hence the \(B_t\)
are i.i.d., with \(\E B_t=1/d\) and \(\operatorname{Var}B_t=\frac{2(d-1)}{d^2(d+2)}\le\frac2{d^2}\). Also
\(\norm{\hat e_T}^2=\prod_t(1-B_t)\le\exp(-\sum_tB_t)\).
Let \(\theta=\eta/4\) and \(L=\Lambda+2\ln\frac1{1-\theta}\). For \(\epsilon\le1/10\) we have \(\Lambda\ge\ln(1/\sin^2 0.1)>4.6\), and
\(2\ln\frac1{1-\theta}\le\frac{8}{3}\theta=\frac23\eta\), so \(L\le\Lambda(1+\eta/6)\). With
\(1+\gamma=(1+\eta)/(1+\eta/6)\), so that \(\gamma\ge5\eta/7\), we get \(T/d\ge(1+\gamma)L\). Chebyshev's inequality gives
\[
 \Prob\Bigl(\sum_{t\le T}B_t<L\Bigr)\le\frac{2T/d^2}{(T/d)^2\gamma^2/(1+\gamma)^2}
 =\frac{2(1+1/\gamma)^2}{T}\le\frac{11.52}{\eta^2T}\le\frac{2.51}{\eta^2d}\le\frac16 ,
\]
using \((1+1/\gamma)\le2.4/\eta\), \(T\ge4.6d\) and \(d\ge16/\eta^2\). Off this event,
\(\norm{\hat e_T}\le e^{-L/2}=(1-\theta)\sin\epsilon\).

\paragraph{Rounding part: contraction in \(L^2\).} The vector \(r_t\) is \(\mathcal F_t\)-measurable, and \(u_{t+1},\dots,u_T\) are
i.i.d.\ uniform and independent of \(\mathcal F_t\). For a fixed \(w\) we have \(\E\norm{P_uw}^2=(1-1/d)\norm w^2\). Iterating,
\(\E[\norm{P_T\cdots P_{t+1}r_t}^2\mid\mathcal F_t]=(1-1/d)^{T-t}\norm{r_t}^2\). By Minkowski's inequality in \(L^2(\Prob)\),
\[
 \bigl(\E\norm{a_T}^2\bigr)^{1/2}\le\sum_{t=1}^T(1-1/d)^{(T-t)/2}\delta_q\le\frac{\delta_q}{1-\sqrt{1-1/d}}\le2d\delta_q=d^{3/2}2^{-q}.
\]
By the choice of \(q\), \(2^q\ge4\sqrt6\,d^{3/2}/(\eta\sin\epsilon)\), so \((\E\norm{a_T}^2)^{1/2}\le\theta\sin\epsilon/\sqrt6\). Markov's inequality on
\(\norm{a_T}^2\) then gives \(\Prob(\norm{a_T}>\theta\sin\epsilon)\le1/6\).

\paragraph{Conclusion.} With probability at least \(1-\frac16-\frac16=\frac23\),
\(\norm{z_T-s}\le\norm{\hat e_T}+\norm{a_T}\le\sin\epsilon<1\), so \(z_T\ne0\). The
point \(z_T\) lies in the closed ball of radius \(\sin\epsilon\) about the unit vector \(s\), and every ray from the origin meeting this ball makes
an angle at most \(\arcsin(\sin\epsilon)=\epsilon\) with \(s\). Hence \(\arccos\ip{z_T/\norm{z_T}}{s}\le\epsilon\). For the bit count,
\(\log_2(1/\sin\epsilon)\le\lambda+0.0025\) for \(\epsilon\le1/10\) (as \(\sin\epsilon\ge\epsilon(1-\epsilon^2/6)\)). Then
\(q+3\le\lambda+\frac32\log_2d+\log_2\frac1\eta+\log_2(4\sqrt6)+4.0025\le\lambda+\frac32\log_2d+\log_2\frac1\eta+7.3\).
Finally, \(0\le\ln\frac1{\sin^2\epsilon}-\ln\frac1{\epsilon^2}=2\ln\frac{\epsilon}{\sin\epsilon}\le-2\ln(1-\epsilon^2/6)\le\epsilon^2/2\), and
\(\frac{\epsilon^2/2}{\ln(1/\epsilon^2)}\le\frac{0.005}{4.605}<0.0011\). \qed

\begin{remark}[Where the bits go]
Rounding cost \(\lambda\) bits plus \(\frac32\log_2 d+O(1)\) bits per coordinate: \(\frac12\log_2d\) for
\(\ell_\infty\to\ell_2\), and \(\log_2 d\) for the effective memory \(\sum_k(1-1/d)^{k/2}\le2d\) of the contraction. The \(L^2\) argument
avoids a \(\log_2 T\approx\log_2(d\lambda)\) term, which is why \(C\) does not depend on \(\epsilon\). We did not try to
optimize the \(\frac32\). The \(\lambda\) term cannot be avoided by any learner: the final output takes at most \(2^M\) values (for each fixed
sample sequence), each successful only for \(\hat s\) in a cap of angular radius \(\epsilon\), whose normalized measure is
at most \(\frac12(\sin\epsilon)^{d-1}\) by the cap bound used in the terminal lemma; averaging over a uniform \(\hat s\), success
probability \(\ge2/3\) forces \(2^M\cdot\frac12(\sin\epsilon)^{d-1}\ge\frac23\), i.e.\ \(M\ge(d-1)\lambda'+\log_2\frac43\). (The last column of
Table~\ref{tab:kacz}, with \(\lceil\lambda'\rceil-3\) fractional bits and no \(\frac32\log_2 d\) term, only illustrates
that the algorithm fails below this count.)
\end{remark}

\begin{remark}[Probability estimate]
Chebyshev's inequality is enough because \(\operatorname{Var}\sum B_t=O(T/d^2)\) while the required deviation is \(\eta T/d\). Bernstein's
inequality for the bounded variables \(B_t\in[0,1]\), or the exact formulas \(\E\ln\frac1{1-B}=\psi(\frac d2)-\psi(\frac{d-1}2)\) and
\(\operatorname{Var}\ln\frac1{1-B}=\psi'(\frac{d-1}2)-\psi'(\frac d2)\), would give exponential tails and a threshold
\(d_0=O(\eta^{-2})\) with a better constant. This is not needed for the first-order statement.
\end{remark}

\section{Where the constant was lost}\label{sec:ledger}

Table~\ref{tab:ledger} evaluates the lower bound \eqref{eq:TLB}, i.e.\ \(T_{\rm LB}/(d\lambda)\), after each successive correction.
Rows S1--S5 use the paper's parameters \(k=m=\lfloor d/16\rfloor\), \(\ell=\lfloor d/32\rfloor\) at \(d=4\cdot10^4\), with exact beta products
(script \texttt{computations/ledger.py}). Row S6 is the parameter limit (Lemma~\ref{lem:sup}).

\begin{table}[h]
\centering\small
\begin{tabular}{@{}lp{8.3cm}ll@{}}
\toprule
step & change & \(c\) (\(\lambda\to\infty\)) & factor\\
\midrule
S0 & paper as printed, \texttt{07-beta-mixture.tex}:117 & \(1.53\times10^{-5}=2^{-16}\) & --\\
S1 & same displayed bounds, no slack (\(k,p\ge d/16\) not \(d/32\); \(\log_2G=4d\) not \(5d\); \(\lambda-3\approx\lambda\)) & \(9.77\times10^{-4}=2^{-10}\) & \(\times64\)\\
S2 & exact values of the stated lemma bounds \(2^{D'}(1+\E Z^{-\ell})\), \(\E\rho^{-p}\) & \(3.88\times10^{-3}\) & \(\times4.0\)\\
S3 & constant the proof actually yields, \(2^{(b-1)/2}(1+\E Z^{-\ell})\) & \(0.1048\) & \(\times27\)\\
S3b & split at \(1/2\), exact beta ratios & \(0.1195\) & \(\times1.14\)\\
S4 & moving split \(q_0=1-\delta\), optimized per \(j\) (Lemma~\ref{lem:moving}) & \(0.8704\) & \(\times7.3\)\\
S5 & terminal cap \(8\to1\) (Lemma~\ref{lem:cap}); at \(\lambda=\log_210\): \(0.0655\to0.8515\) & \(0.8704\) & (\(\times13\) at \(\lambda=\log_2 10\))\\
S6 & parameters \(\mu\ll\kappa\ll\pi\to0\) (Lemma~\ref{lem:sup}) & \(\to2\ln2=1.3863\) & \(\times1.59\)\\
 & certified uniform constant at \((10^{-7},10^{-4},10^{-8})\) & \(1.38470\) & total \(\times9.07\cdot10^4\)\\
\bottomrule
\end{tabular}
\caption{Ledger of the constant. The largest single losses are the arithmetic slack (S1, \(\times64\)), the
statement-versus-proof gap of the mixture lemma (S3, \(\times27\)), and the fixed split (S4, \(\times7.3\)). The terminal
constant matters only for bounded \(\lambda\), where it decides whether the bound is uniform at all.}
\label{tab:ledger}
\end{table}

\subsection{Limit of the original framework}\label{sec:framework}
Keep the original split \(q_0=1/2\) but free all parameters \((\kappa,\pi,\mu)\) and use terminal constant \(1\). The first
component of the mixture then costs \(\frac{b-1}2\ln2\) per direction, which is \(\frac\kappa2\ln2\) nats per unit dimension per block.
This holds whatever the parameters are. With the paper's ratio estimate \(\B(\alpha,\beta_1)\le\B(\alpha,\beta_0)\), the second
component additionally costs \(\E Z^{-\ell}=e^{(\kappa/2)d\ln\frac1{1-\pi}(1+o(1))}\). As \(\kappa,\mu\to0\) the rate becomes
\(g_{1/2}=\frac\kappa2[\ln2+2\ln\frac1{1-\pi}]\), and
\[
 c_{1/2}=\sup_{0<\pi<1}\frac{2\pi\ln2}{\ln2+2\ln\frac1{1-\pi}}=\frac{\ln2}{u}=0.33357\ldots,\qquad
 u=-W_{-1}\!\Bigl(-\frac{1}{e\sqrt2}\Bigr)=2.07796\ldots,
\]
attained at \(\pi^*=1-1/u=0.51876\). Here \(W_{-1}\) is the lower branch of the Lambert function. The first-order condition
\(\ln2+2\ln u=2(u-1)\), with \(u=1/(1-\pi)\), gives \(c_{1/2}=\ln2\,(1-\pi^*)\). Numerical optimization over all
\((\kappa,\pi,\mu)\) in the original framework confirms \(0.333571\) (\texttt{computations/rates.py}, mode \texttt{half}), which is
\(4.16\) times below \(2\ln 2\). With the exact beta ratio but still \(q_0=1/2\), the second component becomes
\(O(1)\), but the first does not. This gives
\(\sup_\pi 2\pi\ln2/\ln\frac2{1-\pi}=2\ln2/u'=0.51759\) with \(u'=-W_{-1}(-1/(2e))=2.67835\). Either way the cap is a
fixed fraction of \(2\ln 2\), so the split point must move to \(1\).

\section{Finite \texorpdfstring{\(d\)}{d}}\label{sec:finite}
The convergence in Theorem~\ref{thm:main} is slow. The constant approaches \(2\ln 2\) only when
\(\mu\ll\kappa\ll\pi\to0\). The integrality requirements \(k,m\ge2\) alone need \(d\ge2/\mu\). The memory term
\(\ln w/m\) must also be small compared with \(d\,g\approx\frac{\kappa}{2}\pi d\), i.e.\ \(M/d^2\ll\mu\kappa\pi\). Table~\ref{tab:finite}
reports the effective constant \(T_{\rm LB}/(d\lambda)\) from~\eqref{eq:TLB} with exact beta products, worst test
dimension \(D_{\min}\), \(\lambda=10\), and Kaczmarz memory \(M=d(\lambda+2\log_2 d)\), maximized over a parameter grid.
These are estimates and not certificates: the block maximum is evaluated at \(b\in\{k-1,k\}\) and \(j\in\{1,m/2,m-1\}\).

\begin{table}[h]
\centering\small
\begin{tabular}{@{}lcccc@{}}
\toprule
\(d\) & \(10^4\) & \(10^6\) & \(10^8\) & \(10^{10}\)\\
\midrule
best value found (estimate, not a certificate) & 0.341 & 0.779 & 1.071 & 1.231\\
\bottomrule
\end{tabular}
\caption{Effective lower-bound constant at finite \(d\), \(\lambda=10\), memory \(d(\lambda+2\log_2d)\) (compare \(2\ln2=1.386\) and the original \(2^{-16}\)).
Each entry is the maximum over the union of two parameter grids (\texttt{computations/independent/indep\_constant.py}, \(\kappa\le0.03,\pi\le0.2\);
\texttt{computations/finite\_d\_fine.py}, \(\kappa\le0.3,\pi\le0.7\)); the \(10^8\) entry comes from the first grid, the others from the second.
Optimizers: \((0.076,0.48,0.13)\), \((0.038,0.33,0.024)\), \((0.01,0.2,0.01)\), \((0.0024,0.10,0.0019)\). Since the block maximum is
evaluated only at \(b\in\{k-1,k\}\), \(j\in\{1,m/2,m-1\}\), these numbers are estimates of what the proof gives, not proved bounds.}
\label{tab:finite}
\end{table}

So at \(d=10^4\) the improved proof gives about \(0.34\) (about \(2.2\cdot10^4\) times the original \(2^{-16}\)), which is still a factor \(4\) below
Kaczmarz's \(1.386\). Even at \(d=10^{10}\) the gap is \(11\%\). The statement
``Kaczmarz is first-order optimal'' is asymptotic. The rate at which the effective constant approaches \(2\ln2\) is
not optimized here.

\section{Numerical unit test of Lemma~\ref{lem:kacz}}
Script \texttt{computations/kacz\_fixedpoint.py} runs exactly the learner of Section~\ref{sec:kaczmarz}, with float64 standing in for
exact arithmetic. The rounding resolution is \(2^{-q}\ge2^{-35}\), far above the float64 error. It uses \(200\) independent
signals per row and \(\eta=1/4\). The lemma's sufficient condition \(d\ge16/\eta^2=256\) fails at \(d=200\). Those rows test the
algorithm, not the lemma's Chebyshev constant.

\begin{table}[h]
\centering\small
\resizebox{\textwidth}{!}{%
\begin{tabular}{@{}rrrrrccccccc@{}}
\toprule
\(d\) & \(\epsilon\) & \(q\) & \(M/d\) & \(T\) & succ.\ fix/real & rms\(\norm{a_T}\) & bound & \(\hat\Prob(\sum B<L)\) & \(\max\norm{\tilde z}_\infty\) & 1st hit & succ.\ \(q_{\rm low}\)\\
\midrule
200 & 0.1 & 21 & 24 & 1153 & 1.000/1.000 & \(2.7\cdot10^{-5}\) & \(1.3\cdot10^{-3}\) & 0.000 & 0.310 & 1.0307 & 0.000 (\(q=1\))\\
200 & 0.001 & 27 & 30 & 3454 & 1.000/1.000 & \(4.3\cdot10^{-7}\) & \(2.1\cdot10^{-5}\) & 0.000 & 0.312 & 0.9953 & 0.000 (\(q=7\))\\
200 & 1e-05 & 34 & 37 & 5757 & 1.000/1.000 & \(3.4\cdot10^{-9}\) & \(1.6\cdot10^{-7}\) & 0.000 & 0.303 & 0.9989 & 0.000 (\(q=14\))\\
400 & 0.1 & 22 & 25 & 2305 & 1.000/1.000 & \(2.8\cdot10^{-5}\) & \(1.9\cdot10^{-3}\) & 0.000 & 0.206 & 1.0036 & 0.000 (\(q=1\))\\
400 & 0.001 & 29 & 32 & 6908 & 1.000/1.000 & \(2.2\cdot10^{-7}\) & \(1.5\cdot10^{-5}\) & 0.000 & 0.216 & 1.0043 & 0.000 (\(q=7\))\\
400 & 0.0001 & 32 & 35 & 9211 & 1.000/1.000 & \(2.7\cdot10^{-8}\) & \(1.9\cdot10^{-6}\) & 0.000 & 0.236 & 1.0043 & 0.000 (\(q=11\))\\
1000 & 0.1 & 24 & 27 & 5761 & 1.000/1.000 & \(1.7\cdot10^{-5}\) & \(1.9\cdot10^{-3}\) & 0.000 & 0.150 & 0.9982 & 0.000 (\(q=1\))\\
1000 & 0.001 & 31 & 34 & 17270 & 1.000/1.000 & \(1.3\cdot10^{-7}\) & \(1.5\cdot10^{-5}\) & 0.000 & 0.152 & 1.0025 & 0.000 (\(q=7\))\\
\bottomrule
\end{tabular}}
\caption{Fixed-point Kaczmarz at the lemma's \(T\) and \(q\). Columns: success rate (fixed point / real
registers); root-mean-square rounding accumulation \(\norm{a_T}\) versus the proved \(L^2\) bound \(d^{3/2}2^{-q}\);
empirical \(\Prob(\sum B_t<L)\); \(\max\norm{\tilde z_t}_\infty\) (the proof needs \(\le3\)); median first time with angular error \(\le\epsilon\),
divided by \(d\ln(1/\sin^2\epsilon)\) (checked every \(50\) steps, so the resolution is \(50/(d\ln(1/\sin^2\epsilon))\le4\%\)); success with only \(\lceil\lambda'\rceil-3\) fractional bits.}
\label{tab:kacz}
\end{table}

All rows succeed with frequency \(1.000\ge2/3\). The rounding term stays \(47\)--\(115\) times below its proved bound, and
\(\norm{\tilde z_t}_\infty\le0.32\). The first hitting time is \((1.00\pm0.03)\,d\ln(1/\sin^2\epsilon)\), which matches the first-order
constant. With \(\lambda-3\) fractional bits the algorithm never succeeds. The unit test does not probe the optimality of
the \(\frac32\log_2d\) term.

\appendix
\section{Line-by-line changes to the original proof}\label{app:changes}
File references are to the \LaTeX{} sources of~\cite{OAIsub} (\texttt{preprints/Subsphere-methods-\ldots/build/}).
\begin{enumerate}[label=C\arabic*.,leftmargin=2.2em]
\item \texttt{main.tex}:13--14 and \texttt{sections/00-introduction.tex}:53--67, 85--88: replace \(2^{-16}\) by
``any \(c<2\ln2\), with threshold \(d_0(c,M)\)''.
\item \texttt{02-subspace-moments.tex}:24--30, 44--56 (eq.~\texttt{s:affine-terminal}): replace
\(\min\{1,(8\epsilon/R)^{D-1}\}\) by \(\min\{1,\frac12(\sin\epsilon/R)^{D-1}\}\), using Lemma~\ref{lem:cap}: centre the chordal ball at
\(\cos\epsilon\,\widehat s\) (radius \(\sin\epsilon\)), do not recentre at a successful
direction, and use the \(\arcsin\) cap.
\item \texttt{07-beta-mixture.tex}:21--48 (Lemma \texttt{s:mixture}): replace \(\Law(v)\le2^{D'}(\nu_0+\nu_1)\) by Lemma~\ref{lem:moving}
with a free \(\delta\in(0,1)\). Proof lines 86--106: split at \(q_0=1-\delta\) and keep the beta ratios.
\item \texttt{07-beta-mixture.tex}:115--118: \(c_*=2^{-16}\) becomes the given \(c<2\ln2\). Choose admissible
\((\kappa,\pi,\mu,\zeta)\) with \(\Phi>c\) (Lemma~\ref{lem:sup}).
\item \texttt{07-beta-mixture.tex}:121--132: keep the short-run reduction with this \(c\). Delete the branch
\(\lambda\le1024\) (lines 130--132) and with it the use of \texttt{03-stopping.tex}:123--135.
\item \texttt{07-beta-mixture.tex}:134--141: parameters \(k=\lfloor\kappa d\rfloor\), \(m=\lceil\mu d\rceil\),
\(\ell=\lfloor\pi d/2\rfloor\), and \(G=w^{1/m}[(1-\delta)^{-k/2}+r_2\delta^{-1/2}]\Ores\) instead of \(2^{4d}w^{1/m}\).
\item \texttt{07-beta-mixture.tex}:149--153: budget \(J\le2c\lambda_s/\kappa+1=o(d)\), and the eventual
conditions~\eqref{eq:eventual} replace \(J\le d/8\), \(m,p\ge d/32\).
\item \texttt{07-beta-mixture.tex}:166--168: invariant \((8\epsilon/r)^p\to(\sin\epsilon/r)^p\). Lines 176--179: base case via
Lemma~\ref{lem:cap}.
\item \texttt{07-beta-mixture.tex}:216: replace \(D\ge3d/4\), \(n\ge11d/16\) by \(D\ge D_{\min}\), \(n-j\ge D_{\min}-k-m\).
\item \texttt{07-beta-mixture.tex}:239--257: replace eq.~\texttt{(s:mixture-inverse)} (\(\E Z^{-\ell}\le2^d\)) and
\(2^dU_0(1+\E Z^{-\ell})\le2^{3d}U_0\) by Lemma~\ref{lem:invmom} and Proposition~\ref{prop:dir}:
conditional factor \(\le[(1-\delta)^{-k/2}+r_2\delta^{-1/2}]U_0\).
\item \texttt{07-beta-mixture.tex}:262--269: \((2^{3d}U_0)^m\to([(1-\delta)^{-k/2}+r_2\delta^{-1/2}]U_0)^m\).
\item \texttt{07-beta-mixture.tex}:271--281: residual moment \(2^\ell\le2^d\) replaced by \(\Ores=e^{(R_1(\kappa,\pi)+o(1))d}\).
This uses monotonicity in \(b\) and \(D\).
\item \texttt{07-beta-mixture.tex}:283--299: replace the arithmetic~\texttt{(s:numeric-arithmetic)} by~\eqref{eq:TLB} and the
limit argument of Section~\ref{sec:proof}.
\item \texttt{07-beta-mixture.tex}:296--298 (side remark): the displayed chain uses \(p\ge d/32\) and \(J\le32c_*\lambda+1\). Both are
looser than needed (\(p\approx d/16\), \(k\approx d/16\)). This is row S1 of Table~\ref{tab:ledger}, and it is not an error.
\end{enumerate}
No step of the original proof was found to be wrong. All changes are sharpenings. The measurability, padding,
Borel-version and short-run arguments (\texttt{01-residual-spheres.tex}, \texttt{03-stopping.tex}:13--88) are used verbatim.

\section{Reproducibility}
Scripts (Python 3, \texttt{numpy}, \texttt{scipy}, \texttt{mpmath}) are in the directory \texttt{computations/} of the
accompanying repository; \texttt{computations/run\_all.sh} runs all of them:
\texttt{ledger.py} (Table~\ref{tab:ledger}),
\texttt{rates.py} (rate functions; framework limit~\(0.333571\)),
\texttt{certify.py} (interval certificate \(1.38470\)),
\texttt{checks.py} and \texttt{independent/indep\_mixture.py} (Monte Carlo of the beta laws, pointwise domination of
Lemma~\ref{lem:moving} on grids including \(q_0\to1\)),
\texttt{inverse\_moment\_check.py} (Lemma~\ref{lem:invmom}, including the counterexample to the tempting bound),
\texttt{independent/indep\_cap\_kacz.py} (Lemma~\ref{lem:cap} for \(D=2,\dots,399\)),
\texttt{independent/indep\_constant.py}, \texttt{finite\_d\_fine.py} (Table~\ref{tab:finite}),
\texttt{finite\_d\_fullscan.py} (full scan over all \((b,j)\) at the optimizers for \(d=10^4,10^6\); it gives the same block maximum as the sampled points),
\texttt{memratio.py} (\(\Theta(d^2)\) remark), and
\texttt{kacz\_fixedpoint.py} (Table~\ref{tab:kacz}).
The scripts in \texttt{computations/independent/} are separate re-implementations that do not import the others.

\section*{Acknowledgement and provenance}
The results in this note were found and checked by Claude (Anthropic) agents working under the direction of
[AUTHOR NAME]. The note has had no human peer review, and nothing in it is formalized in a proof assistant. The underlying architecture and all unchanged lemmas are from the OpenAI
preprint~\cite{OAIsub}.

\begin{thebibliography}{99}
\bibitem{OAIsub} OpenAI. \emph{Subsphere methods for memory--sample lower bounds in noiseless Gaussian regression.}
OpenAI Math Release preprint, September 27, 2026. \url{https://github.com/openai/math}
\bibitem{OAImem} OpenAI. \emph{Memory and precision in noiseless Gaussian regression.}
OpenAI Math Release preprint, September 27, 2026. \url{https://github.com/openai/math}
\bibitem{SSV19} V.~Sharan, A.~Sidford, G.~Valiant. Memory-sample tradeoffs for linear regression with small error.
\emph{STOC 2019}, 890--901. doi:10.1145/3313276.3316403, arXiv:1904.08544.
\bibitem{MSSV22} A.~Marsden, V.~Sharan, A.~Sidford, G.~Valiant. Efficient convex optimization requires superlinear memory.
\emph{COLT 2022}, PMLR 178:2390--2430.
\bibitem{BZJ23} M.~Blanchard, J.~Zhang, P.~Jaillet. Quadratic memory is necessary for optimal query complexity in convex
optimization: center-of-mass is Pareto-optimal. \emph{COLT 2023}, PMLR 195:4696--4736.
\bibitem{Kac37} S.~Kaczmarz. Angen\"aherte Aufl\"osung von Systemen linearer Gleichungen.
\emph{Bull.\ Int.\ Acad.\ Polon.\ Sci.\ Lett.\ A} 35 (1937), 355--357.
\bibitem{SV09} T.~Strohmer, R.~Vershynin. A randomized Kaczmarz algorithm with exponential convergence.
\emph{J.\ Fourier Anal.\ Appl.} 15(2) (2009), 262--278.
\end{thebibliography}
\end{document}
